\begin{figure*}
\begin{promptbox}{Example 1: ``Language'' Counterfactual Instruction Prompt}
Your task is to summarize a passage in a way which would receive a score following a set criteria.

You will be given the rubric, the desired scores for each part of the rubric, a passage to summarize, and a reference summary which received a different score.
Use this information to write a summary which would correspond to the given scores and while keeping your summary as close to the reference as possible.

\#\#\# Rubric:

\#\#\# Language Beyond Source Text

**Scores:**

- 1: Poor

- 2: Fair

- 3: Good

- 4: Excellent

**Score Descriptions:**

- **1**: Summary shows a very basic understanding of lexical and syntactic structures.

- **2**: Summary shows an understanding of lexical and syntactic structures.

- **3**: Summary shows an appropriate range of lexical and syntactic structures.

- **4**: Summary shows an excellent range of lexical and syntactic structures.

\#\#\# Passage:

\{\{ reading\_passage\_text \}\}

\#\#\# Desired scores:

**Language Beyond Source Text:** \{\{ language\_score \}\}

\#\#\# Reference Summary:

\{\{ original\_student\_summary \}\}

End your response directly after giving the summary

Do not provide explanation, justification, or notes before or after the summary.

Do not preface with "Here is the summary..."

Do not postface the summary with ``Note'':

\#\#\# Summary:
\end{promptbox}
\caption{Example instruction prompt for ``Language'' criterion including reference. This prompt structure is the basis for Gumbel Machine (GM) and Minimal Edit (ME) structures. In-context Learning (ICL) additionally includes examples after the score descriptions (not shown).}
\label{apx:prompt_ref}
\end{figure*}
